\section{近期研究（2026 年 8--11 月）}
\paragraph{HiCo-GS：让八叉树锚点共享层级上下文并约束几何}
2026.08~\cite{2608_14136}。八叉树多层原语能应对城市视距变化，但不同层的锚点各自优化，会割裂粗结构与细节。HiCo-GS 在上层、当前层和下层间查询并融合上下文，再约束由高斯协方差得到的法向与由深度估计的法向接近。边缘感知平滑用来维持大片表面，同时避免把明显图像边界抹平。

论文在 Mill19、UrbanScene3D、MatrixCity 和 China-Pagoda 比较渲染和几何结果。几何约束利用城市表面近似平面的性质，在强曲率、非平面和纹理与几何边界不一致的区域，应单独观察其作用。

\paragraph{ABCD：用预合成的前景和背景训练当前空间分块}
2026.08~\cite{2608_27735}。ABCD 每轮只优化一个空间分块。对相机和活动凸分块，暂时冻结的场景按前后关系预合成为图像，当前分块的渲染再与它们合成。优化因此能保留来自冻结部分的遮挡与外观，而更新显存主要放活动参数和合成缓冲。

论文在 NeRF 的 kitchen 与 garden 比较常规 3DGS、消融合成的分块方法和完整方法。作者指出样本规模小，光栅、优化器和图像缓冲在显存中占比不小。文中的渐近显存推断与这两个场景实际测得的节省应分开看。

\paragraph{Atlas：为头显双目和磁盘流式调度共同设计高斯渲染}
2026.09~\cite{2609_02352}。Atlas 将多数城市高斯放在磁盘，按头显当前所需内容异步调入，利用帧间连贯性减少细节层级搜索，并复用双眼渲染的部分计算。这里优化的是大场景的运行时存储和显示路径，而非从图像学习高斯。

论文对软件双目渲染做图像质量比较，并用 RTL 实现与移动 GPU 和专用渲染硬件比较。双目协议使用固定眼间距和分辨率。专用架构的速度与能耗结果依赖模拟或硬件配置，不能直接视为任意现有头显都可达到的实测值。

\paragraph{InceptionGS：将场景几何作为生成补全的条件}
2026.09~\cite{2609_02747}。大场景的随手采样会造成视角密度高度不均，常规高斯拟合在弱观测区域容易留下细节空缺。InceptionGS 先优化基础 3DGS，同时把场景特定的层次几何用作生成先验的条件，再以光度混合和随视角变化的采样将生成信号逐步反馈到高斯。

论文在 GigaNVS 与 Mip-NeRF 360 等场景比较重建和生成基线。生成细节应与真实图像可验证的内容区分。作者在结论中还列出两个阶段均需逐场景计算，限制其即时处理大批量场景。

\paragraph{OceanXL：按点预算分块并轻量密度控制}
2026.09~\cite{2609_29985}。大范围水下视频重建既要控制单块优化规模，又要保留远处结构。OceanXL 用二元空间划分使各块点数符合 GPU 预算，并在块内采用水下适配的密度控制与选择性高斯差分激活。划分边界随点分布调整，而不是始终等距切分。

原文在 ISRO 分块例中将预算设为每块 4 万点，并和固定中点的 BlockGaussian 切分作对比；这个预算是实验配置。多水下数据对照给出画质与计算成本折中。作者说明当前深度分层固定，也未显式处理随时间变化的场景。

\paragraph{Budgeted-GS：以容量预算约束训练和细节层级渲染}
2026.10~\cite{2610_03162}。Budgeted-GS 一方面将已有的高斯模型整理成可按预算调度的细节层级树，另一方面在训练时按估计的容量要求设定原语上限，在增密过程中控制模型大小。两条路径分别服务已有模型的部署和已知预算下的直接训练。

作者报告训练侧匹配预算的比较以及大场景渲染实验。论文限于单 GPU 设置，透明场景、极端遮挡和跨存储流式渲染尚未验证。容量分析只应按文中定义的渲染器和一致性泛函理解。

\paragraph{RoMo-GS：在线建图中对曝光运动做解析积分}
2026.10~\cite{2610_24605}。RoMo-GS 把曝光期间投影高斯沿近似线性轨迹的贡献写成解析积分，用一次光栅化逼近以往多时刻采样平均的结果。在线后端联合修正位姿、速度和高斯，另以局部优化与基于 Hessian 的评分抑制原语增长。

作者在合成模糊、公开视觉数据和机器人现场录制中比较图像质量和速度。当前积分依赖曝光时长和曝光期线性运动假设，轨迹弯曲、相机参数不明或前端跟踪失败均会影响映射质量。

\paragraph{ReScene：把稀疏复访中的高斯变化写成可回放记录}
2026.11~\cite{2611_20433}。复访通常只提供少量新图像，又未必能取回早期帧。ReScene 根据图像差异在高斯合成时的贡献，把变化归因给部分原语，用增补和可回放的更新记录保留过去版本。已经保存的版本可以按记录恢复，这里的精确指对历史模型的复现。

论文分开比较当前重建、历史回放、几何与存储，并以 Multi-State、WAT 和 ChronoScene 为测试。方法要求提供统一坐标中的标定位姿，复访间缺少连续观测。它能保留各次重建，但不能反推中间的真实运动，也会保存早期模型已有的错误。

\paragraph{DOOM-GS：以时空分块和流式泼溅训练城市级动态高斯}
2026.11~\cite{2611_22805}。动态场景离开显存训练时，一个视角的可见原语也可能塞不进显存。DOOM-GS 按空间与时间组织数据块，利用相邻相机可见集合重叠来排程，并以流式泼溅和梯度估计在有界显存下更新模型。

论文在加入程序生成车辆和行人的 DynamicMatrixCity 比较单一动态模型与逐帧静态基线，并使用留出相机质量、静态区闪烁和训练时间衡量。结果针对该模拟城市场景和相机划分，真实交通动态的效果需另行验证。
